\section{Bose--Einstein and Fermi--Dirac distributions}
\label{v:rec:ocupacion:section:02}

The equilibrium formulation starts from the Hamiltonian and the total
occupation number,
\begin{equation}
 H=\sum_kE_k^{\rm occ} n_k,
 \qquad
 N=\sum_k n_k,
 \qquad
 K_\mu:=H-\mu N.
 \label{v:rec:eq:hamiltoniano-gran-canonico-rev3}
\end{equation}
For \(\beta>0\), the grand canonical state is
\begin{equation}
 \rho_{\beta,\mu}
 =\frac{e^{-\beta K_\mu}}{\mathcal Z_{\beta,\mu}},
 \qquad
 \mathcal Z_{\beta,\mu}
 =\operatorname{Tr}(e^{-\beta K_\mu}).
 \label{v:rec:eq:estado-gran-canonico-rev3}
\end{equation}

\begin{theorem}[Grand canonical factorization and occupation laws]
\label{v:rec:thm:factorizacion-gran-canonica-rev3}
For energy levels \(E_k^{\rm occ}\) with degeneracies \(g_k\), the
fermionic trace factorizes as
\begin{equation}
 \mathcal Z_{\rm FD}
 =\prod_k\bigl(1+e^{-\beta(E_k^{\rm occ}-\mu)}\bigr)^{g_k},
 \qquad
 \langle n_k\rangle_{\rm FD}
 =\frac{g_k}{e^{\beta(E_k^{\rm occ}-\mu)}+1}.
 \label{v:rec:eq:factorizacion-fd-rev3}
\end{equation}
In the bosonic sector, under the condition
\(\mu<\inf_kE_k^{\rm occ}\),
\begin{equation}
 \mathcal Z_{\rm BE}
 =\prod_k\bigl(1-e^{-\beta(E_k^{\rm occ}-\mu)}\bigr)^{-g_k},
 \qquad
 \langle n_k\rangle_{\rm BE}
 =\frac{g_k}{e^{\beta(E_k^{\rm occ}-\mu)}-1}.
 \label{v:rec:eq:factorizacion-be-rev3}
\end{equation}
\end{theorem}

\begin{proof}
The operator \(K_\mu\) is a sum of commuting occupation operators, so
its trace factorizes over modes. In each fermionic mode, the
occupations \(0\) and \(1\) are summed, and the corresponding factor is
\(1+u_k\), where
\(u_k=e^{-\beta(E_k^{\rm occ}-\mu)}\). In each bosonic mode, all
nonnegative occupations are summed, giving the geometric series
\((1-u_k)^{-1}\); the condition on \(\mu\) ensures \(u_k<1\). The
powers \(g_k\) incorporate the degeneracy. Finally,
\[
 \langle n_k\rangle
 =-\frac1\beta\frac{\partial}{\partial E_k^{\rm occ}}
   \log\mathcal Z
\]
yields the two occupation laws
\cite{Bose1924,Einstein1924,Fermi1926,Dirac1926,Pathria2011}.
\end{proof}
